% LLM-assisted: moved from the longer Methods during the length audit.
% LLM-assisted: Codex drafted these procedures from code and saved runs,
% 4 November 2026. Earlier source credits are retained.
% This records the implemented study, not student review or new experiments.
\section{Full experiment specifications}
\label{app:experiment-settings}
\subsection{Data and preprocessing}
The Runge experiments shared one
800-point development set and 200 points of reserved test data.
The classification experiments kept publications together when splitting development
and test images. Within its development data, separate monitoring sets
controlled early stopping and validation folds ranked hyperparameters.

\subsubsection{Runge regression}
We generated 1000 evenly spaced inputs on $[-1,1]$ and targets
\[
y_i=\frac{1}{1+25x_i^2}+\epsilon_i,
\qquad \epsilon_i\sim\mathcal N(0,0.1^2),
\]
with NumPy's \nolinkurl{linspace} with both endpoints included. Gaussian
noise was drawn with \nolinkurl{default_rng(42)}. We saved the exact Runge
values alongside the noisy targets as a reference for the plots.
We trained the Runge models on the noisy targets and measured their
MSE against those targets.

Scikit-learn's \nolinkurl{train_test_split}, with
\nolinkurl{test_size=0.2}, \nolinkurl{shuffle=True} and
\nolinkurl{random_state=42}, split the data into 800 development points
and 200 test points, without stratification. We applied
\nolinkurl{KFold(n_splits=5, shuffle=True, random_state=42)} to the
development set. Each cross-validation split had 640 training points and 160 validation
points.
We used the same training and validation points for every Runge model
comparison.

Following Section~\ref{sec:theory-scaling}, we fitted the input scaling
separately within each cross-validation split. We standardised each input
feature with the training points' mean and population standard deviation (\nolinkurl{ddof=0}). We also
subtracted the training points' target mean from $y$, without scaling its
variance. We transformed the validation data with the statistics from the training set. We then added the training target mean back to predictions to return
them to the original scale.

For the Runge task, each network received the standardised value of $x$
as its only input feature. For the polynomial models, we formed the features
$x,x^2,\ldots,x^p$ and then standardised each feature column. Centring the
targets handled the polynomial intercept, so we did not add a constant
feature column. If a feature had zero standard deviation, we used one as
its scaling divisor. Before the final fits, we recomputed the preprocessing
statistics from all 800 development points and used them to transform
the test data.

The 200 test points had already been scored in Project 1 and earlier
Runge runs. We reused these points as the reserved test data for the current comparisons. The current searches picked hyperparameter configurations from validation scores.
The random split shows how well the models predict within $[-1,1]$.

\subsubsection{Greek letter images}
We downloaded \nolinkurl{ALPUB_v2.zip} on 27 October 2026 from the author's
\href{https://www.kaggle.com/datasets/miswindall/al-pub-v2}{Kaggle release},
linked from the \href{https://data.cs.mtsu.edu/al-pub/}{MTSU dataset page}.
The dataset's background is described by
\textcite{swindall2021greek}. Our audit found 205797 JPEG images, each
$70\times70$ pixels, in 24 supplied letter classes.
Lunate sigma is a C-shaped form of the Greek letter sigma, common in
the ancient papyri covered by the dataset. In AL-PUB v2, it replaced
the earlier release's sparsely represented sigma class. This change kept
the task at 24 classes and reduced the severe shortage of sigma images
\parencite[Sec.~7.3]{swindall2024papyrology}.
The software successfully read all JPEG files and converted them into pixel values. Two sets of identical pixels carried conflicting
labels: five wholly white and sixteen wholly black images. We excluded
all 21, leaving 205776 images, and kept the raw archive unchanged.
The audit checked archive integrity and exact pixel duplicates.

We used Pillow to place each image on a white background and convert it
to grayscale, then applied \nolinkurl{ImageOps.autocontrast}. We resized the images with Lanczos resampling to fit inside $26\times26$ pixels while
keeping their aspect ratio. We centred each image on a white $28\times28$ canvas and flattened it
row by row into 784 pixel features for the network. Pixel intensity $p$ was
mapped to $1-p/255$, so white became zero and darker ink approached one.
We kept grayscale values without binary thresholding. The exact Pillow
calls and centring rule are given in Appendix~\ref{app:methods-reproduction}.

We used no augmentation or class weighting. These per-image transforms
needed no fitted statistics across the dataset. We stored the processed pixel values as unsigned 8-bit integers and
converted them to float64 when training or assessing a model. We kept the
same output order for all 24 letter classes. For training, we one-hot encoded each image's target: a vector with one
at its class position and zero elsewhere.

The final experiments reused the publication-based holdout assignments
generated during our project's initial AL-PUB data checks. We grouped images by the POxy publication item
number in each filename. Images from the same publication may share visual
characteristics. We kept them together to reduce the risk of overly
optimistic test results and to test how well the models work on other
publications.

We reserved 284 of the 1419 publication groups for testing. This left
159053 development images from 1135 groups and 46723 test images from
284 groups. Publication groups do not guarantee different writers in the
development and test sets. The grouping rule and exact split procedure
are recorded in Appendix~\ref{app:methods-reproduction}.

\subsection{Shared network and training rules}
Our NumPy FFNN implemented the row-batch forward pass in
Eq.~\eqref{eq:forward}. It stored the layer values needed by the manual
backward sweep in Eqs.~\eqref{eq:backward}--\eqref{eq:parameter-gradients}.
The Runge networks had one linear output and full MSE from
Eq.~\eqref{eq:mse}. The letter networks had 24 softmax outputs and
mean cross-entropy, with the output derivative in
Eq.~\eqref{eq:ce-gradient}. We used the actual number of samples in each
batch for normalisation and applied no second batch average.
Cross-entropy was evaluated from logits with log-sum-exp, as described
in Section~\ref{sec:theory-losses}. JAX supplied independent reference
gradients for checks and was not used to train these networks.

We initialised each network's biases at zero. We initialised the weights
of sigmoid hidden layers and all output layers with Xavier normal draws with variance $2/(n_{\rm in}+n_{\rm out})$.
For ReLU and Leaky ReLU hidden layers, we used He normal draws with variance
$2/n_{\rm in}$ \parencite{hjorthjensenLectureNotes}. Leaky ReLU had fixed
negative slope 0.01. At zero, the implementation used derivative zero
for ReLU and 0.01 for Leaky ReLU. We kept separate optimiser state for each layer's weight matrix and
bias vector. For each fit, we reset a NumPy \nolinkurl{default_rng}
for initialisation and a separate one for shuffling, with the seeds
stated for that experiment. We drew weights with \nolinkurl{Generator.normal}
and shuffled the training samples each epoch with \nolinkurl{Generator.permutation}.
The shuffle generator continued from its state in the preceding epoch. If the last minibatch was smaller, we computed its gradients using the
actual number of samples in that minibatch. We fixed the RMSProp decay at 0.9.
For Adam, we fixed $\beta_1=0.9$ and $\beta_2=0.999$ and applied bias
correction. Both optimisers had $\varepsilon=10^{-8}$. The base learning rate stayed constant within each fit and was picked as set out below.
All network computations ran in float64.

Regularised networks used the weight-only L1 and L2 cost functions of
Section~\ref{sec:theory-regularisation}. The penalty gradients were added
to the data-loss gradients before the Adam moment updates. Biases were
excluded, and the penalty had no further division by sample or parameter
count. The L1 subgradient at zero was set to zero. We recorded the
penalised training cost and used unpenalised prediction scores for
comparisons across penalties \parencite{goodfellow2016deep}.

\subsection{Runge experiments}
\subsubsection{Optimiser comparison}
The training rules of Section~\ref{sec:theory-training} can reach different
fits under the same epoch budget. We tested their achieved MSE and
fitting cost while holding the hidden activation and initialisation
scheme fixed. We compared a network with one hidden layer of 50 neurons with a network
that had two hidden layers of 100 neurons each. Both had sigmoid hidden
activation functions and no penalty. For each architecture we tested full-batch gradient descent (GD), minibatch RMSProp and minibatch Adam at base learning rates
$\{0.001,0.01,0.03,0.1\}$. All fits ran for 1000 epochs. GD used all
training points per update, while RMSProp and Adam used batches of 32.
Initialisation and shuffle seeds were fixed at 2026. At the fold-training size of 640, each minibatch epoch made 20 updates,
while each GD epoch made one. The comparison therefore varied the
update rule together with the number of updates per epoch. It assessed
these training setups under a common pass budget.

We fitted each of the 24 hyperparameter configurations on all five cross-validation splits.
For each architecture and optimiser, we picked the learning rate that gave
the lowest mean validation MSE. We then fitted the networks with those
six picked hyperparameter configurations on all 800 development points. We also picked the best of those six configurations by mean validation
MSE. OLS was fitted by pseudoinverse
at fixed degrees 14 and 15 on the same folds and development rows, with
\nolinkurl{numpy.linalg.pinv} with its default cutoff.
Degree 14 came from the earlier validation-based selection. Degree 15 came from an earlier
test-score sweep. We kept it as a historical comparison with that
test-guided choice of degree recorded.

\subsubsection{Activation and network size}
We compared training and validation MSE as width and depth changed. We next compared sigmoid, ReLU and Leaky ReLU, one or two hidden layers,
and widths 25, 50 and 100 in each hidden layer. We ran Adam because it gave the lowest mean validation MSE in the
optimiser comparison in Section~\ref{sec:results-optimisers}. For each activation--architecture
pair, we tested learning rates $\{0.0003,0.001,0.003\}$ and picked the learning rate
that gave the lowest mean validation MSE. We assessed 54 hyperparameter configurations on the same five cross-validation splits.
Ten of the needed fits matched completed sigmoid/Adam fits from the
optimiser comparison, so we reused their saved weights and scores.
We trained the other 260 fits anew. Each network was trained for 1000
epochs with batches of 32 and seed 2026. Its output layer was linear,
and no penalty was applied.

Each activation--architecture pair had its own validation-selected
learning rate in the network-size comparison. We computed validation MSE
minus training MSE within each fold, then averaged those paired gaps.
We also recorded whether a larger model lowered training error while
raising validation error, using the overfitting pattern in
Section~\ref{sec:theory-regularisation}. Changing the hidden activation function also changed its weight-initialisation
scheme under the rules above. The comparison therefore includes that
change. We used these validation folds to pick the learning rates, so
their scores describe performance on the validation sets used for hyperparameter selection.
No test scores or explicit bias--variance decomposition were produced in
this experiment.

\subsubsection{Regularisation and polynomial regression models}
\label{app:regularisation-search}
We compared prediction errors across the tuned unregularised, L1 and L2
FFNN families. These manual
FFNNs used sigmoid hidden activation functions and Adam, with hidden widths
$[25]$, $[25,25]$ or $[50,50]$, informed by the capacity comparison. We kept
1000 epochs and batches of 32. Scikit-learn's \nolinkurl{ParameterSampler},
with \nolinkurl{n_iter=16} and \nolinkurl{random_state=42}, drew combinations
from the architecture list in the order shown above. Learning rates used SciPy's
\nolinkurl{loguniform} distribution on $[10^{-4},10^{-2}]$, and positive
regularisation parameters used the same distribution on $[10^{-7},10^{-2}]$.
The L1 and L2 searches used the same sampled hyperparameter configurations.
For each combination, we also fitted a network with the same architecture
and learning rate but no penalty, giving 16 unregularised comparison hyperparameter configurations. Random search gave us a bounded
search budget \parencite{bergstra2012random}.

We fitted every sampled configuration on the five shared folds at
seed 2026 and picked the lowest mean unpenalised validation MSE within
each family. Each selected network was fitted on all 800 development
points at seed 2026 and scored on the shared 200 test points.

Ridge used degree 14 and $\lambda=10^{-6}$, the five-fold pick in
Project 1's \nolinkurl{results/part_d_best_by_k.csv}, with historical
validation MSE 0.0118443. Lasso used degree 6 and $\lambda=10^{-4}$,
the lowest five-fold validation MSE among candidates marked converged
in every fold in \nolinkurl{results/part_i_common_cv.csv}, at 0.0180245.
The later Ridge grid omitted $10^{-6}$, and the earlier Lasso test
minimum was unconverged. The ten-fold Ridge pick used a different
assessment protocol. We kept the five-fold selections and fitted them
anew on the shared folds and development set. The corrected Ridge baseline used the normal
equations with $n\lambda I$, matching full MSE, solved with
\nolinkurl{numpy.linalg.solve}. We fitted Lasso with the reused full-batch Adam/JAX gradient solver with learning rate 0.01 and at most 10000 updates. Its initial
coefficients were drawn from $\mathcal N(0,0.1^2)$ with a fresh
\nolinkurl{default_rng(42)}. It stopped when the Euclidean norm of the
change in coefficients fell below $10^{-7}$.

We checked Ridge against augmented least squares, solved with
\nolinkurl{numpy.linalg.lstsq} and \nolinkurl{rcond=None}. The Lasso reference
was scikit-learn's \nolinkurl{Lasso}, with $\alpha=\lambda/2$ to match its
half-MSE convention, \nolinkurl{fit_intercept=False},
\nolinkurl{tol=1e-12} and \nolinkurl{max_iter=100000}. Its coordinate selection
used the default cyclic order. Convergence flags and cost function differences
were saved.

\subsection{Letter-classification experiment}
We applied the multiclass loss of Section~\ref{sec:theory-losses} to the
784 pixel features. The model comparison tested whether learned hidden
features improved macro-F1 over the no-hidden-layer softmax baseline.
A matched library FFNN checked performance under a separately implemented
training path. We used a fixed search subset of 24000 development images to keep the search
within the available runtime. Scikit-learn's
\nolinkurl{StratifiedShuffleSplit(n_splits=1, train_size=24000, random_state=2026)}
picked the search subset while keeping roughly the original letter proportions.

We assessed each hyperparameter configuration with scikit-learn's
\nolinkurl{StratifiedGroupKFold(n_splits=5, shuffle=True, random_state=2026)}.
It balanced letter proportions as far as the publication groups allowed,
while keeping each publication in one fold. We used the validation set to measure each model's predictive
performance. The images left after setting aside a validation fold formed
that split's outer training pool.

Within each outer training pool, we used a second
\nolinkurl{StratifiedGroupKFold} with five folds and shuffling enabled.
Its seed was $2027+f$, where $f=0,\ldots,4$ labels the outer
cross-validation split. We used
its first split to define the training set and monitoring set. The monitoring set controlled early stopping, keeping the validation
fold separate from the stopping decisions as set out in
Section~\ref{sec:theory-assessment}. No publication appeared in more than one of the training, monitoring or
validation sets within a cross-validation split. Each set contained images
from all 24 letter classes.

After each training epoch, we measured cross-entropy on the monitoring
set. We stopped after ten epochs without a strict decrease in that loss,
or after 100 epochs, and restored the weights from the epoch with the
lowest monitoring cross-entropy. We did not use validation scores to decide
when to stop. All classification models used batches of 128 and seed 2026.
We recorded whether each fit stopped early or reached the epoch ceiling.
Convergence remained unverified.

For the minibatch and epoch-monitoring workflow, we consulted the public
NumPy example in \textcite[Ch.~11]{raschka2022machine}. Our network kept
its own softmax cross-entropy and grouped monitoring.

For the classification learning curves, training CE was averaged over
minibatches before each update. Successive batches therefore used
changing weights. The training and validation scores in the result
tables were instead computed at one fixed restored model per fold.

Scikit-learn's \nolinkurl{ParameterSampler(n_iter=12, random_state=2026)}
drew the hyperparameter configurations. We tested one to three hidden layers with equal
widths picked from $\{32,64,128\}$. Activations were sigmoid, ReLU or
Leaky ReLU. Adam learning rates followed SciPy's \nolinkurl{loguniform} distribution
on $[10^{-4},10^{-2}]$. Penalties were none, L1 or L2, with positive
strengths drawn from \nolinkurl{loguniform} on $[10^{-7},10^{-2}]$.
Unregularised hyperparameter configurations had regularisation parameter zero.
Each hyperparameter configuration had the same five
cross-validation splits.

We ranked hyperparameter configurations by mean validation macro-F1.
Macro-F1 weights letters equally despite their unequal class sizes. We also saved
accuracy and mean cross-entropy. The largest softmax output gave the
predicted class. Macro-F1 included all 24 labels, with undefined per-class
F1 set to zero. Exact sampling order and metric calls are recorded in
Appendix~\ref{app:methods-reproduction}.

Logistic regression used the manual softmax implementation with no hidden
layer or penalty. It tested learning rates $\{10^{-4},10^{-3},10^{-2}\}$ on the same
folds and stopping scheme. The scikit-learn FFNN
\parencite{pedregosa2011sklearn} matched the highest-scoring sampled
no-penalty sigmoid/ReLU hyperparameter configuration in architecture, activation, learning rate,
batch size and precision. We used \nolinkurl{MLPClassifier} with Adam
and no penalty, training one epoch at a time through \nolinkurl{partial_fit}.
This let the scikit-learn FFNN continue training from its preceding epoch
while our grouped monitoring rule controlled early stopping. Each fit used NumPy's \nolinkurl{RandomState(2026)}.
The scikit-learn FFNN kept its own initialisation and Adam implementation.
This was a matched predictive comparison, with those implementation
differences still present. The exact
constructor settings are in Appendix~\ref{app:methods-reproduction}.
We tested 12 hyperparameter configurations for the manual FFNN and three
learning rates for logistic regression. We fitted the scikit-learn FFNN with
one matched hyperparameter configuration.

Before assessing test performance, we fixed the picked hyperparameter configurations and final
training budgets. The final fits used all development images, so each
epoch processed more images than a search-fold epoch. To carry the
monitoring-based stopping evidence into a bounded final fit, we used the
training-image exposure measure introduced in
Section~\ref{sec:letter-methods}. We set the final
budget to match the median number of training-image exposures at the best
monitoring epochs, rounding up to a whole epoch and using at least one:
\[
E_{\rm final}=\max\!\left(1,
\left\lceil\frac{\operatorname{median}_{f}
(E_{{\rm best},f}n_{{\rm fit},f})}{159053}\right\rceil\right).
\]
Here $E_{{\rm best},f}$ is the best monitoring epoch and
$n_{{\rm fit},f}$ is the number of training images in fold $f$.
We fitted the final classification models on every development image,
without holding any back for monitoring. The epoch budget followed
the exposure-matching rule above.
We then assessed the three fixed models on the same 46723 test images.
Test scores did not guide model settings or epoch budgets for these
classification experiments.
We saved confusion counts and per-class precision, recall and F1 alongside
the overall scores. 

\subsection{Checks, uncertainty and reproducibility}
We compared manual backpropagation against independent float64 JAX
cost functions for the losses, activations, depths and penalties used here
\parencite{bradbury2018jax}. Checks included zero weights, nonzero biases,
activation behaviour at zero, extreme logits, seeded initialisation and
short final minibatches. A linear network on matching polynomial features
was checked against OLS. We also tested early stopping and weight
restoration. Separate inference routines reloaded saved weights and
recomputed scores and predictions. We checked that publication groups
stayed separate across the classification data partitions and that each
training, monitoring and validation set contained all 24 classes. The elementwise pass rule and gradient tolerances for each experiment are
listed in Appendix~\ref{app:check-tolerances}. Numerical outcomes appear
at the start of Results, in Section~\ref{sec:implementation-results}.

For each hyperparameter configuration, we reported the mean validation
score and its sample standard deviation across five folds, with divisor
four. Fold SD measures variation across validation folds. All reported
network comparisons use training seed 2026.

We saved the software versions with each experiment, alongside the
split membership and sampled hyperparameter configurations. Saved weights and predictions
give a record of the assessed models. Appendix~\ref{app:methods-reproduction}
gives the software versions and run entry points. It also identifies the
saved files needed to recover the exact splits and draws.
Original source archives and dated snapshots stay unchanged.

